\section{限界と今後の課題}
\label{sec:limitations}

\paragraph{コートの汎化と幾何教師}
学習規模を拡張して視点分布の広さと密度を独立に制御し、未知の収録で低角・横向き・見切れ別に採点する。
合成データの投影教師にはSfMの時間依存の誤差と地面の凹凸が入る。再訪時の幾何整合と独立の地面基準から両者を調べる。
実写の固定カメラでも、振動・ズーム・レンズ歪みや同期の残差は三角測量に伝わる。

\paragraph{IDと絶対3D精度}
停止クリップの介入では、修正したboxからのpose・外観の再抽出と、人物対応通過後の全時系列のID採点が残る。
少数の選手から作った疑似参照は、採点に使った観測と由来が重なる。全画面の人物検出の適合率は未集計である。
独立参照に対する球位置・選手関節・rootの絶対3D誤差は未計測である。
モーションキャプチャや独立に校正した観測を同じ時刻とコート座標に揃え、IDの一致、位置誤差、欠測を同時に測る。

\paragraph{球像と時系列}
初版の外部注釈は生成モデルと実行IDが未確認である。
母集団の画面外率と遮蔽率を測るには、層化した目視例に加え、蓄積した推定座標を照合する独立の参照が要る。
単フレームと時系列の差にはMDDの寄与が含まれる。時間長の寄与を調べるには、入力特徴を揃えた比較を複数seedと独立testに広げる。

\paragraph{実行系の測定範囲}
横断バッチの効果を測るには、クリップ規模、形状分布、保存済み成果物、OSのキャッシュ状態を揃えた反復計測が要る。
横断バッチを外したTensorRT構成の計時が残る。実行条件間の出力差は、座標差の最大値と有効性マスクの変化も追う。
